% ===================================================================
%  TABLICA Nr 5 — Dom i jego budowa.
%  Meme regime que les precedents.
%
%  CE TABLEAU EST UN DIALOGUE D'UN BOUT A L'AUTRE, et les attributions
%  de parole se composent « \textsc{...}. --- », comme dans l'ido : le
%  releve en compte trente-six, et c'est a ce nombre que kolonoj.py
%  reconnait un dialogue (voir « narracio »). Le polonais n'a pas de
%  niveaux de langue grammaticalises comme le javanais ; mais le neveu
%  y vouvoie les deux hommes de metier et tutoie son oncle, et cette
%  difference-la est dans le texte, non dans une regle.
%
%  LES NOMS PROPRES DE L'IDO SONT DES MOTS : « siorulo Blanko » et
%  « siorulo Rufo » sont M. Blanc et M. Roux. On les traduit — pan
%  Biały, pan Rudy — comme les prenoms, et non on les translittere.
%
%  LE VOCABULAIRE DU BATIMENT EST LE PLUS DENSE DU LIVRET : cent
%  cinquante renvois pour un seul tableau. Chaque metier a son nom
%  polonais propre — murarz, cieśla, dekarz, tynkarz, stolarz, malarz,
%  szklarz, ślusarz, kowal — et aucun n'est un mot general.
% ===================================================================

%%K t05-apar-1 apar
\VUsaut{6.0mm}
\VUtitre{15.0pt}{60}{TABLICA Nr 5}
\VUsaut{1.4mm}
\VUtitre{11.4pt}{0}{Dom i jego budowa.}
\VUsaut{2.2mm}

%%K t05-tit-1 sub
\VUpk{10.2pt}{40}{Szczęśliwe spotkanie.}
\VUsaut{3.0mm}

%%K t05-01-1 p
\textsc{Jan}. --- Wuju, bardzo chciałbym wiedzieć, jak się buduje \VUgras{dom}.

%%K t05-01-2 p
\textsc{Wuj} (80). --- To bardzo łatwe, mój przyjacielu, chodź ze mną, pokażę ci ten,
który pan Bernard (79) każe budować i w którym będę mieszkał.

%%K t05-01-3 p
\textsc{Jan}. --- Kto więc narysował \VUgras{plan} \textsuperscript{(76)} tego domu?

%%K t05-01-4 p
\textsc{Wuj}. --- \VUgras{Architekt} \textsuperscript{(75)}; przedstawił bardzo wyraźny
rysunek, który został zatwierdzony, i \VUgras{przedsiębiorca} \textsuperscript{(77)}
zaczął roboty.

%%K t05-01-5 p
\textsc{Jan}. --- A kto jest przedsiębiorcą?

%%K t05-01-6 p
\textsc{Wuj}. --- To pan Biały, ojciec twojego przyjaciela Jakuba. Kieruje robotnikami
razem z panem Rudym, architektem.

%%K t05-01-7 p
No! Oto w samą porę idzie pan Biały ze swoją \VUgras{linijką} \textsuperscript{(78)}, a
pan Rudy ze swoim planem.

%%K t05-01-8 p
Dzień dobry panom. Jak się panowie miewają?

%%K t05-01-9 p
\textsc{Pan Rudy}. --- Bardzo dobrze, dziękuję. Dzień dobry, Janie, jak urosło to
dziecko!

%%K t05-01-10 p
\textsc{Wuj}. --- Tak, doprawdy, to mały mężczyzna. Jest bardzo poważny, trzeba go
objaśniać o wszystkim, co widzi. Dzisiaj chciałby wiedzieć, jak się buduje dom.

%%K t05-tit-2 sub
\VUpk{10.2pt}{40}{Robotnicy.}
\VUsaut{3.0mm}

%%K t05-01-11 p
\textsc{Pan Biały}. --- No, chodź ze mną, mój mały, zaraz ci to wytłumaczę, bo bardzo
lubię dzieci, które chcą się uczyć.

%%K t05-01-12 p
Widzisz tego robotnika, który trzyma \VUgras{łopatę} \textsuperscript{(82)} w ręku: to
\VUgras{kopacz} \textsuperscript{(81)}. Kopie \VUgras{rów} \textsuperscript{(46)}, żeby
położyć rury wodociągowe. Wykopał też ziemię w miejscu, gdzie stoi dom, żeby murarz
wzniósł cztery \VUgras{mury}. Część tych murów, która jest w ziemi, nazywa się
\VUgras{fundamentami} \textsuperscript{(2)}. Przestrzeń zawarta między fundamentami
tworzy \VUgras{piwnicę} \textsuperscript{(3)}. Ta piwnica ma małe okno zwane
\VUgras{okienkiem piwnicznym} \textsuperscript{(5)}, \VUgras{drzwi}
\textsuperscript{(4)} i schodki.

%%K t05-01-13 p
\VUgras{Murarz} \textsuperscript{(108)} dalej budował mury, zostawiając otwory na
\VUgras{drzwi} \textsuperscript{(8)} i \VUgras{okna} \textsuperscript{(17)}.

%%K t05-01-14 p
\textsc{Jan}. --- Ale tego murarza nie widzę!

%%K t05-01-15 p
P. \textsc{Biały}. --- Teraz skończył pracować w domu. Jest obok \VUgras{stajni}
\textsuperscript{(54)}, na swoim \VUgras{rusztowaniu} \textsuperscript{(110)}, buduje
mur \VUgras{pralni} \textsuperscript{(56)}.

%%K t05-01-16 p
Ma kielnię, koryto i \VUgras{pion} \textsuperscript{(109)}. Ale nie zapamiętasz tych
nazw.

%%K t05-01-17 p
\textsc{Jan}. --- Owszem, bo zaraz poproszę wuja o małą kielnię i koryto, a sam zrobię
pion ze sznurka.

%%K t05-tit-3 sub
\VUsaut{3.0mm}
\VUpk{10.2pt}{40}{Materiał.}
\VUsaut{3.0mm}

%%K t05-01-18 p
\textsc{Wuj}. --- Ha! bardzo dobrze. Ale powiedz mi, Janie, czego trzeba, żeby zbudować
dom?

%%K t05-01-19 p
\textsc{Jan}. --- Trzeba \VUgras{ciosów} \textsuperscript{(59)} i \VUgras{zaprawy}
\textsuperscript{(65)}.

%%K t05-01-20 p
\textsc{Wuj}. --- Słusznie, zobacz tego człowieka w wielkim kapeluszu, na swoim
\VUgras{wozie} \textsuperscript{(136)}. To \VUgras{woźnica} \textsuperscript{(135)}.
Codziennie przywozi tu \VUgras{bloki kamienne} \textsuperscript{(60)}. Wyładowuje swój
wóz na podwórzu, a \VUgras{kamieniarze} \textsuperscript{(83)} obciosują bloki i tną je
swoją \VUgras{piłą do kamienia} \textsuperscript{(85)}. \VUgras{Robotnik}
\textsuperscript{(146)} nosi je do murarza, który je układa.

%%K t05-01-21 p
Tymczasem \VUgras{zarabiacz} \textsuperscript{(87)} przygotowuje zaprawę. Potrzebuje
\VUgras{piasku} \textsuperscript{(62)}, \VUgras{wapna} \textsuperscript{(63)} i
\VUgras{wody} \textsuperscript{(64)}. Wapno jest obok niego w \VUgras{worku}
\textsuperscript{(71)}, \VUgras{pomocnik murarza} \textsuperscript{(111)} przywozi mu
piasek na \VUgras{taczce} \textsuperscript{(112)}, a \VUgras{uczeń}
\textsuperscript{(113)} stoi przy pompie (58). Napełnia \VUgras{wiadro}
\textsuperscript{(114)}. Zaprawa będzie wkrótce gotowa. \VUgras{Noszący zaprawę}
\textsuperscript{(107)} bierze ją i niesie w górę po drabinie na rusztowanie, gdzie
pracuje murarz, który kończy tę stronę domu.

%%K t05-01-22 p
A teraz jak się nazywa robotnik, który robi \VUgras{dach} \textsuperscript{(31)}?

%%K t05-01-23 p
\textsc{Jan}. --- To \VUgras{dekarz} \textsuperscript{(118)}.

%%K t05-tit-4 sub
\VUsaut{3.0mm}
\VUpk{10.2pt}{40}{Dach domu.}
\VUsaut{3.0mm}

%%K t05-01-24 p
P. \textsc{Biały}. --- Tak, ale najpierw przychodzi \VUgras{cieśla}
\textsuperscript{(115)}.

%%K t05-01-25 p
Cieśla robi \VUgras{więźbę} \textsuperscript{(35)}. Składa \VUgras{belki}
\textsuperscript{(37)} \VUgras{kołkami} \textsuperscript{(117)}. Ci robotnicy, tam w
górze, w swoich szerokich aksamitnych spodniach, to cieśle. Jeden trzyma belkę, a drugi
przycina kołek swoją \VUgras{siekierą} \textsuperscript{(116)}. Ten, który jest obok
\VUgras{komina} \textsuperscript{(41)}, to dekarz. Przybija łaty na (*) legarach i
\VUgras{łupek} \textsuperscript{(66)} na łatach. Czasem kładzie dachówki.

%%K t05-noto-1 noto
\VUnotes{143.2mm}{%
(*) \VUgras{Balk-o} = niem. \textit{Balken}, ang. \textit{joist}, fr. \textit{solive},
wł. \textit{travicello}, ros. \textit{balka}, hiszp. i port. \textit{viga}. Rdzeń
techniczny dotąd nieprzyjęty, ale godny zaproponowania dla oddania znaczenia słowa fr.
\textit{solive}, które nie jest ani belką fr. \textit{poutre}, ani beleczką. Jest to
okrągławo obrobiony drąg drewniany o przekroju kwadratowym, który podtrzymuje podłogę i
sufit.%
}

%%K t05-01-26 p
Dach stajni jest \VUgras{kryty dachówką} \textsuperscript{(55)}, dach domu jest
\VUgras{kryty łupkiem} \textsuperscript{(32)}, a dach werandy jest \VUgras{cynkowy}
\textsuperscript{(24)}.

%%K t05-01-27 p
\textsc{Jan}. --- Czy dekarz robi też dachy cynkowe?

%%K t05-01-28 p
P. \textsc{Biały}. --- Nie, ale \VUgras{blacharz} \textsuperscript{(121)} albo
\VUgras{ołowiarz}, ten, który zakłada \VUgras{świetlik} \textsuperscript{(38)} na dachu
domu. Zakłada też \VUgras{rury spustowe} \textsuperscript{(43)} i \VUgras{rynny}
\textsuperscript{(42)}. Lutuje cynk i blachę. Umocował \VUgras{piorunochron}
\textsuperscript{(40)} i \VUgras{chorągiewkę} \textsuperscript{(39)} na
\VUgras{szczycie} \textsuperscript{(36)}.

%%K t05-01-29 p
\textsc{Jan}. --- Więc dom jest ukończony?

%%K t05-tit-5 sub
\VUsaut{3.0mm}
\VUpk{10.2pt}{40}{Narzędzia.}
\VUsaut{3.0mm}

%%K t05-01-30 p
P. \textsc{Biały}. --- Nie całkiem. \VUgras{Tynkarz} \textsuperscript{(99)} buduje z
\VUgras{cegieł} \textsuperscript{(61)} ścianki, które dzielą go na rozmaite pokoje.
Powleka \VUgras{gipsem} \textsuperscript{(70)} \VUgras{sufity} \textsuperscript{(26)} i
mury. Teraz robi sufit \VUgras{werandy} \textsuperscript{(33)}. Ma, tak jak murarz,
\VUgras{kielnię} \textsuperscript{(100)} i \VUgras{koryto} \textsuperscript{(101)}.

%%K t05-01-31 p
Potem stolarz robi drzwi, okna, \VUgras{parkiety} \textsuperscript{(16)} i
\VUgras{okiennice} \textsuperscript{(19)}.

%%K t05-01-32 p
Teraz pracuje na podwórzu przy swoim \VUgras{warsztacie} \textsuperscript{(90)}. Ma
\VUgras{piłę} \textsuperscript{(94)} do cięcia drewna (69), \VUgras{strug}
\textsuperscript{(91)}, \VUgras{ścisk} \textsuperscript{(92)}, \VUgras{dłuto}
\textsuperscript{(94 bis)}, \VUgras{cyrkiel} \textsuperscript{(95)} i drewniany
\VUgras{młotek} \textsuperscript{(95 bis)}.

%%K t05-01-33 p
\textsc{Jan}. --- Ha! ale zabawniej jest stolarzyć niż murować. Wuju, kupisz mi
warsztat, piłę i strug, bo wkrótce zrobię budę dla Medora.

%%K t05-01-34 p
\textsc{Wuj}. --- Tak, chłopcze.

%%K t05-01-35 p
\textsc{Jan}. --- Jakże jestem zadowolony! A kto to jest ten, który stoi obok drzwi
wejściowych?

%%K t05-tit-6 sub
\VUsaut{3.0mm}
\VUpk{10.2pt}{40}{Malowanie.}
\VUsaut{3.0mm}

%%K t05-01-36 p
P. \textsc{Biały}. --- To \VUgras{malarz} \textsuperscript{(102)}. Stoi na swojej
drabinie z garnkiem \VUgras{farby} \textsuperscript{(103)} i swoim \VUgras{pędzlem}
\textsuperscript{(104)}. Maluje drewno drzwi wejściowych, żeby się dłużej zachowało.

%%K t05-01-37 p
On też klei papiery w pokojach.

%%K t05-01-38 p
\VUgras{Tapicer} \textsuperscript{(107)}, który jest na \VUgras{drabinie}
\textsuperscript{(106)}, na \VUgras{balkonie} \textsuperscript{(28)}, zakłada
\VUgras{roletę} \textsuperscript{(29)}.

%%K t05-01-39 p
Robotnik stojący obok wielkiego okna to \VUgras{szklarz} \textsuperscript{(96)}.
Umocowuje \VUgras{szyby} \textsuperscript{(18)} \VUgras{kitem} \textsuperscript{(73)}.
Ten drugi robotnik, klęczący obok drzwi piwnicy, to \VUgras{ślusarz}
\textsuperscript{(97)}. Zakłada \VUgras{zamek} \textsuperscript{(6)}. Jego
\VUgras{pilnik} \textsuperscript{(98)} leży obok niego.

%%K t05-01-40 p
\textsc{Jan}. --- Czy ten robotnik, który uderza swoim \VUgras{młotem}
\textsuperscript{(84)} w żelastwo, to też ślusarz?

%%K t05-01-41 p
P. \textsc{Biały}. --- Ściślej mówiąc, to kowal (125). Kuje \VUgras{żelaziwa}
\textsuperscript{(67)} dla \VUgras{elektryka} \textsuperscript{(124)}, który jest na
dachu \VUgras{wozowni} \textsuperscript{(51)}. Ma \VUgras{kuźnię}
\textsuperscript{(126)}, \VUgras{kowadło} \textsuperscript{(127)} i \VUgras{obcęgi}
\textsuperscript{(128)}.

%%K t05-01-42 p
Elektryk umocowuje \VUgras{miedziany} \textsuperscript{(44)} \VUgras{drut} telefonu.

%%K t05-tit-7 sub
\VUpk{10.2pt}{40}{Uprawa ogrodu.}
\VUsaut{3.0mm}

%%K t05-01-43 p
\textsc{Wuj}. --- W głębi \VUgras{podwórza} \textsuperscript{(45)}, obok \VUgras{kraty}
\textsuperscript{(47)} i \VUgras{bramy} \textsuperscript{(48)} widzisz
\VUgras{ogrodnika} \textsuperscript{(129)}. Przygotowuje \VUgras{ogródek}
\textsuperscript{(49)}. Przekopał ziemię swoim \VUgras{szpadlem} \textsuperscript{(131)},
sieje nasiona i grabi \VUgras{grabiami} \textsuperscript{(130)}. Potem swoją
\VUgras{konewką} \textsuperscript{(132)} podleje \VUgras{grządki}
\textsuperscript{(150)} i rośliny będą rosły.

%%K t05-01-44 p
\VUgras{Stajenny} \textsuperscript{(133)}, który właśnie oporządził konia i dał mu
paszę, czyści \VUgras{powóz} \textsuperscript{(52)} gąbką, \VUgras{pompką ręczną}
\textsuperscript{(134)} i wiadrem wody. Potem wprowadzi go z powrotem do wozowni i
zamknie drzwi grubym \VUgras{ryglem} \textsuperscript{(53)}.

%%K t05-01-45 p
No i jesteś zadowolony, jak sądzę, miałeś dosyć objaśnień.

%%K t05-01-46 p
\textsc{Jan}. --- Tak, wuju, ale chętnie zobaczyłbym wnętrze domu.

%%K t05-01-47 p
\textsc{Wuj}. --- To będzie jutro. Pan Rudy wytłumaczy ci plan i wskaże ci nazwę każdego
pokoju.

%%K t05-tit-8 sub
\VUsaut{3.0mm}
\VUpk{10.2pt}{40}{Wewnętrzny układ domu.}
\VUsaut{2.4mm}

%%K t05-apar-2 apar
{\centering\textit{(Zobacz plan.)}\par}
\VUsaut{2.4mm}

%%K t05-01-48 p
\textsc{Architekt}. --- Chodź, Janie, zaraz każę ci zwiedzić rozmaite części domu.

%%K t05-01-49 p
Wejdźmy na \VUgras{parter} \textsuperscript{(7)}. Oto guzik \VUgras{dzwonka}
\textsuperscript{(11)}. Wchodzimy po trzech \VUgras{stopniach} \textsuperscript{(14)}
\VUgras{ganku} \textsuperscript{(9)}, przekraczamy \VUgras{próg} \textsuperscript{(10)}
\VUgras{drzwi} wejściowych \textsuperscript{(8)} i oto jesteśmy u stóp \VUgras{schodów}
\textsuperscript{(13)}.

%%K t05-01-50 p
\textsc{Jan}. --- To jest korytarz.

%%K t05-01-51 p
\textsc{Architekt}. --- Nie, to jest \VUgras{przedsionek} \textsuperscript{(12)}.
\VUgras{Korytarz} \textsuperscript{(a)} jest na końcu. Na prawo oto \VUgras{salon}
\textsuperscript{(b)} i \VUgras{jadalnia} \textsuperscript{(c)}, naprzeciw jadalni jest
\VUgras{kuchnia} \textsuperscript{(d)} i \VUgras{pokój służbowy} \textsuperscript{(e)},
potem z przodu \VUgras{gabinet} \textsuperscript{(f)}. \VUgras{Weranda}
\textsuperscript{(23)} ze swoimi \VUgras{oszklonymi drzwiami} \textsuperscript{(25)}
jest obok salonu.

%%K t05-01-52 p
Zaraz wejdziemy po schodach. Trzymaj się \VUgras{poręczy} \textsuperscript{(15)} i licz
stopnie. Jest ich czternaście. Oto \VUgras{sień} \textsuperscript{(m)}, jesteśmy na
pierwszym \VUgras{piętrze} \textsuperscript{(27)}.

%%K t05-tit-9 sub
\VUsaut{3.0mm}
\VUpk{10.2pt}{40}{Pierwsze piętro.}
\VUsaut{3.0mm}

%%K t05-01-53 p
Naprzeciw schodów jest \VUgras{ustęp} \textsuperscript{(20)} i \VUgras{garderoba}
\textsuperscript{(j)}. Na prawo piękna \VUgras{sypialnia} \textsuperscript{(g)} z dwoma
oknami na \VUgras{fasadzie} \textsuperscript{(1)} domu. Na lewo jest \VUgras{łazienka}
\textsuperscript{(1)} i \VUgras{pokój gościnny} \textsuperscript{(i)} z wielką
\VUgras{alkową} \textsuperscript{(h)}. Obok sypialni jest \VUgras{pokój dziecinny}
\textsuperscript{(k)}. Znajduje się przy obszernym \VUgras{tarasie}
\textsuperscript{(21)}. Pod tym tarasem jest drugi ustęp (20), a na murze oto
\VUgras{dzwon} \textsuperscript{(22)}, który dzwoni na obiad.

%%K t05-01-54 p
\textsc{Jan}. --- Chodźmy na \VUgras{drugie piętro}.

%%K t05-01-55 p
\textsc{Architekt}. --- Nie ma drugiego piętra. Pod dachem są tylko \VUgras{mansardy}
\textsuperscript{(33)} i strych (34). Skończyliśmy zwiedzanie, możemy zejść z powrotem.

%%K t05-01-56 p
\textsc{Jan}. --- Dziękuję panu, to bardzo ciekawe. Zaraz opowiem ojcu wszystko, co
widziałem.
\VUsaut{4.0mm}
\VUfilet{22mm}
